% !PRINTED 605
\weberpage{605}
%<<CONT
trary degree $n$, with real coefficients, which have two conjugate imaginary and $n - 2$ real roots; for in
\[
f(x) = (x - \alpha_1)(x - \alpha_2) \ldots (x - \alpha_n)
\]
\[
= x^n + a_1 x^{n-1} + a_2 x^{n-2} + \cdots + a_n
\]
one may take $\alpha_1, \alpha_2$ conjugate imaginary at will and the remaining $\alpha$ real. But the number of real and imaginary roots of $f(x)$ does not change when the coefficients are varied continuously within certain finite limits (cf. the seventh section), and consequently there are also such equations whose coefficients are rational numbers, since one can always find rational numbers in arbitrary nearness to any given real numbers.

It remains, then, only to show that these rational coefficients can be so chosen that $f(x)$ becomes irreducible. But this follows from the following theorem:

3. \gesperrt{If $p$ is a prime, $c_0, c_1, c_2 \ldots c_n$ a sequence of integers, of which $c_0$ is not divisible by $p$, while $c_1, c_2 \ldots c_n$ are divisible by $p$, but $c_n$ is not divisible by $p^2$, then the function}
\[
\varphi(x) = c_0 x^n + c_1 x^{n-1} + \cdots + c_n
\]
\gesperrt{is irreducible.}

For if $\varphi(x)$ splits into two factors with rational coefficients, then (by \S.\,2) the factors may also be assumed to have integral coefficients. So let
\[
\varphi(x) = (\alpha_0 x^h + \alpha_1 x^{h-1} + \cdots + \alpha_h)(\beta_0 x^k + \beta_1 x^{k-1} + \cdots + \beta_k),
\]
$h$ and $k$ both greater than zero and their sum equal to $n$.

Since $\alpha_h \beta_k = c_n$, one of the two factors $\alpha_h, \beta_k$ must be divisible by $p$, and the other not. Suppose that $\beta_k$ is divisible by $p$ and $\alpha_h$ not divisible; and since not all the $\beta$ can be divisible by $p$, because otherwise $c_0$ too would be divisible by $p$, let $\beta_\nu$ be not divisible by $p$, and $\beta_{\nu+1}, \beta_{\nu+2} \ldots \beta_k$ divisible by $p$. The coefficient of $x^{k-\nu}$ in the product of the two factors is then $\alpha_h \beta_\nu + \alpha_{h-1}\beta_{\nu+1} + \cdots$, and hence not divisible by $p$. It would therefore have to be that $k - \nu = n$, which is not possible, since already $k < n$ must hold.

If, then, one sets for the coefficients of $f(x)$:
\[
a_1 = \frac{c_1}{c_0}, \quad a_2 = \frac{c_2}{c_0} \cdots a_n = \frac{c_n}{c_0},
\]
% [page alignment mirrored 2026-08-07] Printed 605 ends immediately after
% the coefficient display above.  The following theorem/prose belongs to
% printed 606 and has been moved to i0631.
%CONT>>
